% Generated draft: requires mathematical and source review.
\documentclass[11pt]{article}
\usepackage[margin=25mm]{geometry}
\usepackage{fontspec}
\setmainfont{Latin Modern Roman}
\usepackage{amsmath,amssymb,mathtools,unicode-math}
\setmathfont{Latin Modern Math}
\usepackage{enumitem,needspace,xurl,hyperref}
\setlength{\parindent}{0pt}
\setlength{\parskip}{5pt}
\setlength{\emergencystretch}{3em}
\newcommand{\N}{\mathbb{N}}
\newcommand{\Z}{\mathbb{Z}}
\newcommand{\Q}{\mathbb{Q}}
\newcommand{\R}{\mathbb{R}}
\newcommand{\C}{\mathbb{C}}
\newcommand{\F}{\mathbb{F}}
\newcommand{\A}{\mathbb{A}}
\newcommand{\PP}{\mathbb{P}}
\newcommand{\E}{\mathbb{E}}
\newcommand{\eps}{\varepsilon}
\newcommand{\dd}{\,\mathrm{d}}
\begin{document}
\section*{UnsolvedMath ID: 2642}
\textbf{Source problem code:} catalogue entry
\subsection{Definitions and mathematical statement}
A subgroup $H$ of a group $G$ is countable here if $|H|\le\aleph_0$, including finite subgroups. If $\mathsf S(H)$ denotes the source's property ``$H$ has a subnormal abelian series,'' its exact local-to-global quantifier pattern is
\[
 \forall G,\qquad
 \bigl(\forall H\le G,\ |H|\le\aleph_0\Longrightarrow\mathsf S(H)\bigr)
 \Longrightarrow\mathsf S(G).
\]
The group $G$ is arbitrary, and the hypothesis allows each countable $H$ to have its own series. No common length, common indexing type, or compatibility of those series is assumed. When $G$ itself is countable the conclusion is already in the hypothesis; the substantive question concerns uncountable groups.

A missing definition prevents replacing $\mathsf S$ by one uniquely justified formal predicate. The supplied record does not say whether ``series'' is finite, an ascending ordinal-indexed series, a descending transfinite series, or a more general linearly ordered subgroup system; it also does not state the limit-stage conditions. For example, one possible ascending convention would require an ordinal $\lambda$ and subgroups $(H_\alpha)_{\alpha\le\lambda}$ with
\[
 H_0=1,\quad H_\lambda=H,\quad
 H_\alpha\triangleleft H_{\alpha+1},\quad
 H_{\alpha+1}/H_\alpha\text{ abelian},
\]
with increasing inclusions and $H_\delta=\bigcup_{\alpha<\delta}H_\alpha$ at each nonzero limit $\delta\le\lambda$. This is an explicit possible convention, not an assertion that it is the one intended by the source. Finite series, descending series with intersection conditions, and general ordered systems must not be silently identified with it. Likewise successive normality $H_\alpha\triangleleft H_{\alpha+1}$ is distinct from requiring every term normal in the whole group. A finite series with successive normality and abelian factors is the familiar solvability condition, and is not automatically the intended infinite-group notion in this question.

Once the correct source definition is recovered, the displayed implication asks whether that precise class is recognizable from all countable subgroups. Every group is the directed union of its countable subgroups, but closure of the required series property under that union is the issue being asked, not an assumption.

The source attributes the problem to M. Trombetti in Kourovka Notebook, New Problems, 21st issue (2026), Problem 21.133. Its literature triage checked 17 August 2026 reports no verified local-to-global theorem or counterexample, labels it open, and suggests seeking a directed-union counterexample or a transfinite extension theorem. Its partial-progress bullet only restates countable locality and supplies no distinct proved result. No source definition or theorem beyond that account is available in the supplied record.
\subsection{Short English statement}
Must a group inherit a subnormal abelian series when every countable subgroup has one? The local-to-global quantifiers are clear, but the source does not specify the kind or indexing of the series, so that definition must be recovered before selecting a precise infinite-group theorem.
\subsection{Sources}
\begin{itemize}
\item Local UnsolvedMath registry, permanent ID $2642$, exact source code \texttt{catalogue entry}, ``Kourovka Notebook Problem 21.133,'' including the literature triage checked 17 August 2026.
\item \emph{Kourovka Notebook}, New Problems, 21st issue (2026), Problem 21.133, attributed to M. Trombetti, \url{https://www.kourovka-notebook.org/}. The supplied record gives no definition page, direct primary-paper citation, or series-convention clarification.
\end{itemize}
\end{document}
